%% toc: P2 11 11.4,11.5,11.6
\section{Duality for arbitrary propositional formulas}
\label{sec:p02-w02-p02-010-duality-arbitrary}

Definition~\ref{def:p02-nnf-dual} defines the dual only on NNF presentations,
and Theorem~\ref{thm:p02-propositional-duality} is stated only for them.  This
section connects that operation with every formula of
\(\mathsf{Fm}_{\mathsf{CPL}}\) through the negation normal form theorem of
Section~\ref{sec:prop-syntax-semantics-04}, which makes every formula
logically equivalent to one in negation normal form.  It then reads the
conjunctive and disjunctive normal forms of
Theorem~\ref{thm:prop-normal-forms} through the dual, and separates
dualization from negation.

\emph{Inputs.}  The language \(\mathsf{Fm}_{\mathsf{CPL}}\) of
Section~\ref{sec:prop-syntax-semantics-01}, with primitives
\(\bottom,\neg,\land,\lor,\to\) and \(\topv\) abbreviating \(\neg\bottom\); the
valuations, truth values \(\mathbf F,\mathbf T\), satisfaction and semantic
notions of Section~\ref{sec:prop-syntax-semantics-02}; the definitions of
literal, clause, \(\operatorname{form}_{\mathsf{CPL}}\), negation normal form,
conjunctive normal form and disjunctive normal form in
Section~\ref{sec:prop-syntax-semantics-04}; the uniform substitutions of
Definition~\ref{def:p02-uniform-substitution}.
\emph{Conclusions.}  Dualization is a well-defined operation on the formulas
in negation normal form
(Lemma~\ref{thm:p02-w02-p02-010-nnf-presentation}); for every formula, all
its NNF duals are logically equivalent to the negation of the formula with
every variable complemented, so they determine one equivalence class, and
dualization reverses consequence and exchanges valid with unsatisfiable
formulas (Theorem~\ref{thm:p02-w02-p02-010-arbitrary-duality}); a conjunctive
normal form of a formula dualizes to a disjunctive normal form of the dual,
and a disjunctive normal form dualizes to a formula equivalent to a
conjunctive normal form of the dual
(Corollary~\ref{thm:p02-w02-p02-010-normal-forms}).  All proofs are in the
annex.

\subsection{Complemented valuations and the complementing substitution}
\label{sec:p02-w02-p02-010-complement}

\begin{hbconvention}{def:p02-w02-p02-010-complement}{Complemented valuation and complementing substitution}
For a valuation \(v\), the \emph{complemented valuation} \(\bar v\) is given by
\(\bar v(p_n)=\mathbf T\) iff \(v(p_n)=\mathbf F\), for every \(n\).  It is the
valuation written \(1-v(p)\) in
Theorem~\ref{thm:p02-propositional-duality}, under the reading
\(\mathbf F=0\), \(\mathbf T=1\).  Complementing twice returns the valuation:
\(\bar{\bar v}=v\); so every valuation \(u\) is \(\bar v\) for exactly one
\(v\), namely \(v=\bar u\).

The \emph{complementing substitution} is the uniform substitution
\(\upsilon_{\neg}\) of Definition~\ref{def:p02-uniform-substitution} with
\(\upsilon_{\neg}(p_n)=\neg p_n\) for every \(n\).  Thus
\(\varphi[\upsilon_{\neg}]\) replaces every variable occurrence of
\(\varphi\) by the negation of that variable, simultaneously, and leaves every
connective in place.  The valuation \(v_{\upsilon_{\neg}}\) of
Definition~\ref{def:p02-uniform-substitution} is \(\bar v\), because
\(\widehat v(\neg p_n)=\mathbf T\) iff \(v(p_n)=\mathbf F\).

A formula in negation normal form is called an \emph{NNF formula}.  If
\(\alpha\) is an NNF formula with \(\alpha\equiv\varphi\), then
\(\alpha^{\mathrm d}\) is called an \emph{NNF dual of} \(\varphi\).
\index{duality!propositional}\index{dual formula}
\end{hbconvention}

The NNF presentations of Definition~\ref{def:p02-nnf-dual} display \(\topv\)
as a separate leaf, while in \(\mathsf{Fm}_{\mathsf{CPL}}\) the string
\(\topv\) is the negation \(\neg\bottom\).  The next lemma shows that this
causes no ambiguity: an NNF formula, read as a string, has only one NNF
presentation, so its dual is determined by the string.

\begin{hblemma}{thm:p02-w02-p02-010-nnf-presentation}{proved=annex:p02-w02-p02-010-nnf-presentation}{Unique NNF presentation}
Every NNF formula \(\alpha\) is exactly one of the following: a variable
\(p_n\); a negated variable \(\neg p_n\); \(\bottom\); \(\topv\), that is
\(\neg\bottom\); a conjunction \((\beta\land\gamma)\); or a disjunction
\((\beta\lor\gamma)\).  In the last two cases \(\beta\) and \(\gamma\) are
uniquely determined and are again NNF formulas.  Consequently the clauses of
Definition~\ref{def:p02-nnf-dual} define \(\alpha^{\mathrm d}\) for every NNF
formula \(\alpha\) by structural recursion, and \(\alpha^{\mathrm d}\) is an
NNF formula.
\end{hblemma}

\subsection{The duality theorem for arbitrary formulas}
\label{sec:p02-w02-p02-010-arbitrary}

\begin{hbtheorem}{thm:p02-w02-p02-010-arbitrary-duality}{proved=annex:p02-w02-p02-010-arbitrary-duality}{Duality for arbitrary formulas}
Let \(\varphi,\psi\in\mathsf{Fm}_{\mathsf{CPL}}\).
\begin{enumerate}
  \item \(\varphi\) has at least one NNF dual.
  \item For every NNF formula \(\alpha\) with \(\alpha\equiv\varphi\) and
        every valuation \(v\),
        \[
         v\entails\alpha^{\mathrm d}
         \quad\text{iff}\quad
         \bar v\not\entails\varphi;
         \qquad\text{equivalently,}\qquad
         \alpha^{\mathrm d}\equiv\neg\bigl(\varphi[\upsilon_{\neg}]\bigr).
        \]
  \item Any two NNF duals of \(\varphi\) are logically equivalent.  If
        \(\varphi\) is itself an NNF formula, \(\varphi^{\mathrm d}\) is one of
        its NNF duals.
  \item Let \(\alpha\) and \(\beta\) be NNF duals of \(\varphi\) and \(\psi\)
        respectively.  Then \(\varphi\entails\psi\) iff
        \(\beta\entails\alpha\), and \(\varphi\equiv\psi\) iff
        \(\alpha\equiv\beta\).
  \item With \(\alpha\) an NNF dual of \(\varphi\): \(\varphi\) is valid iff
        \(\alpha\) is unsatisfiable, and \(\varphi\) is unsatisfiable iff
        \(\alpha\) is valid.
\end{enumerate}
\end{hbtheorem}

Clause~3 is the precise sense in which ``the dual of \(\varphi\)'' is
determined: not as a string, since different NNF formulas equivalent to
\(\varphi\) have different duals, but up to logical equivalence.  Clause~4
reverses the direction of consequence; clause~5 transfers the
valid--unsatisfiable exchange of Theorem~\ref{thm:p02-propositional-duality}
from NNF presentations to all formulas.

\begin{hbcorollary}{thm:p02-w02-p02-010-normal-forms}{proved=annex:p02-w02-p02-010-normal-forms}{Normal forms of the dual}
Let \(\varphi\in\mathsf{Fm}_{\mathsf{CPL}}\).
\begin{enumerate}
  \item If \(\gamma\) is a formula in conjunctive normal form with
        \(\gamma\equiv\varphi\), then \(\gamma^{\mathrm d}\) is an NNF dual of
        \(\varphi\) and is in disjunctive normal form.  Such a \(\gamma\)
        exists by Theorem~\ref{thm:prop-normal-forms}.
  \item If \(\delta=\bigvee_{i=1}^n\bigwedge_{j=1}^{k_i}L_{ij}\) is a formula
        in disjunctive normal form with \(\delta\equiv\varphi\), then
        \(\delta^{\mathrm d}\) is an NNF dual of \(\varphi\), and it is
        logically equivalent to the formula in conjunctive normal form built
        from the clauses \(\{L_{i1},\ldots,L_{ik_i}\}\), \(1\leq i\leq n\)
        (the empty conjunction \(\topv\) when \(n=0\)).  It need not itself be
        in conjunctive normal form: \((p_1\land p_1)\) dualizes to
        \((p_1\lor p_1)\).
\end{enumerate}
\end{hbcorollary}

\subsection{Worked example}
\label{sec:p02-w02-p02-010-example}

\begin{hbexample}{def:p02-w02-p02-010-example}{Computing an NNF dual through a conjunctive normal form}
Let \(\varphi=((p_1\to p_2)\land\neg(p_2\land\neg p_3))\).  The code order of
the literals involved is \(\neg p_1<p_2<\neg p_2<p_3\) (codes \(3,4,5,6\)).
Eliminating \(\to\) and moving the negation inward, as in the negation normal
form theorem, gives
\[
 \gamma=((\neg p_1\lor p_2)\land(\neg p_2\lor p_3)),
\]
which is the conjunctive normal form built from the clauses
\(\{\neg p_1,p_2\}\) and \(\{\neg p_2,p_3\}\), and \(\gamma\equiv\varphi\).  Its
dual is
\[
 \gamma^{\mathrm d}=((\neg p_1\land p_2)\lor(\neg p_2\land p_3)),
\]
a disjunctive normal form and, by
Corollary~\ref{thm:p02-w02-p02-010-normal-forms}(1), an NNF dual of
\(\varphi\).  Check of
Theorem~\ref{thm:p02-w02-p02-010-arbitrary-duality}(2) at the valuation \(v\)
with \(v(p_1)=v(p_2)=\mathbf T\), \(v(p_3)=\mathbf F\): both disjuncts of
\(\gamma^{\mathrm d}\) contain a false literal (\(\neg p_1\) and \(\neg p_2\)
respectively), so \(v\not\entails\gamma^{\mathrm d}\).  The complemented
valuation has \(\bar v(p_1)=\bar v(p_2)=\mathbf F\) and
\(\bar v(p_3)=\mathbf T\); then \(p_1\to p_2\) is true, \(p_2\land\neg p_3\) is
false, and so \(\bar v\entails\varphi\).  Both sides of the identity fail, as
they must.  The other form of clause~2 gives the same verdict:
\(\varphi[\upsilon_{\neg}]=((\neg p_1\to\neg p_2)\land
\neg(\neg p_2\land\neg\neg p_3))\) is true under \(v\), so its negation is false.

For the valid--unsatisfiable exchange, \((p_1\lor\neg p_1)\) is a valid NNF
formula and its dual \((p_1\land\neg p_1)\) is unsatisfiable.
\end{hbexample}

\subsection{Dualization is not negation}
\label{sec:p02-w02-p02-010-not-negation}

\begin{hbwarning}{def:p02-w02-p02-010-not-negation}{Dual, negation, and connective swapping}
Three operations are easily confused.
\begin{enumerate}
  \item \emph{The dual is not the negation.}  The literal \(p_1\) is its own
        dual, while its negation is \(\neg p_1\); under any \(v\) with
        \(v(p_1)=\mathbf T\), \(p_1\) and \(\neg p_1\) differ.  By
        Theorem~\ref{thm:p02-w02-p02-010-arbitrary-duality}(2) the dual is
        logically equivalent to the negation of
        \(\varphi[\upsilon_{\neg}]\), not of \(\varphi\): negation alone
        evaluates \(\varphi\) at \(v\), dualization evaluates it at \(\bar v\)
        and then negates.  The dual of an NNF formula is again an NNF formula
        with the same literals, whereas the negation of an NNF formula whose
        main connective is \(\land\) or \(\lor\) is not in negation normal
        form.
  \item \emph{Swapping connectives in a formula that is not in negation normal
        form does not compute a dual.}  In \((p_1\to p_2)\) there is no
        \(\land\) or \(\lor\) to swap, so the swap returns \((p_1\to p_2)\).
        Its NNF dual is \((\neg p_1\land p_2)\), the dual of
        \((\neg p_1\lor p_2)\).  The valuation with
        \(v(p_1)=v(p_2)=\mathbf F\) satisfies \((p_1\to p_2)\) but not
        \((\neg p_1\land p_2)\), so the two are not equivalent.  The negation
        normal form must be formed first.
  \item \emph{Duality is not closed under substitution of schematic letters.}
        The equivalence \((p_1\lor(p_2\land\neg p_2))\equiv p_1\) holds, and
        its dual equivalence \((p_1\land(p_2\lor\neg p_2))\equiv p_1\) holds by
        Theorem~\ref{thm:p02-propositional-duality}.  But if the first
        equivalence is read as an instance of the schema
        \(A\lor B\equiv A\) with \(A=p_1\) and \(B=(p_2\land\neg p_2)\), the
        schema obtained by swapping only the displayed connective,
        \(A\land B\equiv A\), has the false instance
        \((p_1\land(p_2\land\neg p_2))\equiv p_1\).  The dual must be computed
        through every connective of the NNF formula, including those inside
        the formulas substituted for schematic letters.
\end{enumerate}
\end{hbwarning}

\subsection{Relationships and sources}
\label{sec:p02-w02-p02-010-relationships}

Theorem~\ref{thm:p02-propositional-duality} and
Table~\ref{tab:p02-duality} give the operation on NNF presentations; the
present section extends their content to \(\mathsf{Fm}_{\mathsf{CPL}}\) through
the negation normal form theorem and Theorem~\ref{thm:prop-normal-forms}.  The
complementing substitution is an instance of
Theorem~\ref{thm:p02-uniform-substitution}.  The conjunctive normal forms used
here preserve truth under every valuation; the satisfiability-preserving
conjunctive normal form with fresh variables is the different construction of
Section~\ref{sec:p02-tseitin-transformation}
(Theorem~\ref{thm:p02-tseitin}, proved in Annex~\ref{annex:p2.prop.tseitin}).

The syntactic form of the duality goes back to Kleene's treatment of the
propositional calculus.  Kleene's Theorem~8 of \S\,27 states that
interchanging conjunction with disjunction and each proposition letter with its
negation, in a formula built from letters and negated letters by conjunction
and disjunction, yields a formula equivalent to the negation of the original
\cite[\S\,27, Theorem~8, p.~121]{Kleene1952}; its corollary, the principle of
duality, transfers an equivalence between two such formulas to their duals
\cite[\S\,27, p.~122]{Kleene1952}, and the second part of that corollary
reverses implications \cite[\S\,27, p.~124]{Kleene1952}.  Kleene states these
results for provable equivalence and implication in his formal system, and
remarks that duality holds there only as a subsidiary deduction rule, not as a
direct rule, and he traces the recognition of duality in logic to Schr\"oder
in 1877 \cite[\S\,27, p.~123]{Kleene1952}.  The statements and proofs in this section
are semantic, for the language and valuations of
Sections~\ref{sec:prop-syntax-semantics-01}
and~\ref{sec:prop-syntax-semantics-02}; clause~3 of the warning above is the
semantic counterpart of Kleene's remark.
